\begin{thebibliography}{54}
\bibitem{ref1} 3GPP, ``Evolved Universal Terrestrial Radio Access (E-UTRA); Radio Resource Control (RRC); Protocol specification,'' 3GPP TS 36.331, Release 17, 2022.

\bibitem{ref7} T. Ankome and E. Hanada, ``Reactive to predictive mobility management: a systematic review of ML-driven handover optimization in 5G and beyond,'' Machine Learning and Knowledge Extraction, vol. 8, no. 5, art. 133, 2026, doi:10.3390/make8050133.

\bibitem{ref8} B. Saoud, I. Shayea, M. A. Alnakhli, and H. Mohamad, ``Mobility and handover management in 5G/6G networks: challenges, innovations, and sustainable solutions,'' Technologies, vol. 13, no. 8, art. 352, 2025, doi:10.3390/technologies13080352.

\bibitem{ref9} C. Chabira, I. Shayea, G. Nurzhaubayeva, L. Aldasheva, D. Yedilkhan, and S. Amanzholova, ``AI-driven handover management and load balancing optimization in ultra-dense 5G/6G cellular networks,'' Technologies, vol. 13, no. 7, art. 276, 2025, doi:10.3390/technologies13070276.

\bibitem{ref54} M. Dzaferagic, B. Missi Xavier, D. Collins, V. D'Onofrio, M. Martinello, and M. Ruffini, ``ML-based handover prediction over a real O-RAN deployment using RAN intelligent controller,'' \textit{IEEE Transactions on Network and Service Management}, vol. 22, no. 1, pp. 635--647, 2025, doi:10.1109/TNSM.2024.3468910.

\bibitem{ref36} S. Kaufman, S. Rosset, C. Perlich, and O. Stitelman, ``Leakage in data mining: formulation, detection, and avoidance,'' ACM Transactions on Knowledge Discovery from Data, vol. 6, no. 4, art. 15, 2012, doi:10.1145/2382577.2382579.

\bibitem{ref56} S. Kapoor and A. Narayanan, ``Leakage and the reproducibility crisis in machine-learning-based science,'' \textit{Patterns}, vol. 4, no. 9, art. 100804, 2023, doi:10.1016/j.patter.2023.100804.

\bibitem{ref3} H. Deng, C. Peng, A. Fida, J. Meng, and Y. C. Hu, ``Mobility support in cellular networks: a measurement study on its configurations and implications,'' in Proc. ACM Internet Measurement Conference (IMC), 2018, doi:10.1145/3278532.3278546.

\bibitem{ref4} M. Ghoshal, I. Khan, P. Dinh, Z. J. Kong, O. Basit, S. Wang, Y. Feng, Y. C. Hu, and D. Koutsonikolas, ``Handover configurations in operational 5G networks: diversity, evolution, and impact on performance,'' arXiv:2511.03116, 2025.

\bibitem{ref5} A. Hassan, A. Narayanan, A. Zhang, W. Ye, R. Zhu, S. Jin, J. Carpenter, Z. M. Mao, F. Qian, and Z.-L. Zhang, ``Vivisecting mobility management in 5G cellular networks,'' in Proc. ACM SIGCOMM, 2022, pp. 86--100, doi:10.1145/3544216.3544217.

\bibitem{ref41} M. Ghoshal, I. Khan, Z. J. Kong, P. Dinh, J. Meng, Y. C. Hu, and D. Koutsonikolas, ``Performance of cellular networks on the wheels,'' in Proc. ACM Internet Measurement Conference (IMC), 2023, doi:10.1145/3618257.3624814.

\bibitem{ref6} Z. Liu, Q. Deng, Z. Tan, Z. Qian, X. Zhang, A. Swami, and S. V. Krishnamurthy, ``M2HO: Mitigating the adverse effects of 5G handovers on TCP,'' in Proc. ACM MobiCom, 2024, doi:10.1145/3636534.3690680.

\bibitem{ref14} D. Zidic, T. Mastelic, I. Nizetic Kosovic, M. Cagalj, and J. Lorincz, ``Analyses of ping-pong handovers in real 4G telecommunication networks,'' Computer Networks, vol. 227, art. 109699, 2023.

\bibitem{ref15} A. Amirova, A. Abdiraman, L. Aldasheva, I. Shayea, D. Yedilkhan, and A. Tussupov, ``Data-driven and machine learning-based analysis of handover behavior and network stability in mobile networks,'' Future Internet, vol. 18, no. 6, art. 290, 2026, doi:10.3390/fi18060290.

\bibitem{ref45} M. Kabeer, R. Nordin, M. Behjati, and S. L. Lau, ``Sequence-based deep learning for handover optimization in dense urban cellular network,'' arXiv:2510.02958, 2025.

\bibitem{ref46} S. V. R. Mandapati, ``Meta-continual mobility forecasting for proactive handover prediction,'' arXiv:2512.11841, 2025.

\bibitem{ref57} P.-H. Chou, D.-C. Lin, H.-Y. Wei, W. Saad, and Y. Tsao, ``Learning-based radio link failure prediction based on measurement dataset in railway environments,'' arXiv:2511.08851, 2026.

\bibitem{ref58} I. Panitsas, A. Mudvari, A. Maatouk, and L. Tassiulas, ``A deep and transfer learning approach for handover management in O-RAN,'' arXiv:2404.08113, 2024 (later versions retitled ``Predictive handover strategy in 6G'').

\bibitem{ref59} K. Boutiba, M. Bagaa, and A. Ksentini, ``Radio link failure prediction in 5G networks,'' in Proc. IEEE Global Communications Conference (GLOBECOM), 2021.

\bibitem{ref60} W.-C. Chien, Y. Huang, B.-Y. Chang, and W.-Y. Hwang, ``Privacy-preserving handover optimization using federated learning and LSTM networks,'' \textit{Sensors}, vol. 24, no. 20, art. 6685, 2024, doi:10.3390/s24206685.

\bibitem{ref61} A. Wadud, F. Golpayegani, and N. Afraz, ``AI-driven multi-modal adaptive handover control optimization for O-RAN,'' arXiv:2603.17158, 2026.

\bibitem{ref62} M. L. A. Bartsioka, A. Giannopoulos, and S. Spantideas, ``Intelligent dynamic handover via AI-assisted signal quality prediction in 6G multi-RAT networks,'' arXiv:2510.14832, 2025.

\bibitem{ref63} K. Sun, Q. Han, Z. Yang, W. Huang, H. Zhang, and V. C. M. Leung, ``Proactive handover type prediction and parameter optimization based on machine learning,'' \textit{IEEE Transactions on Wireless Communications}, vol. 24, no. 4, pp. 3515--3528, 2025, doi:10.1109/TWC.2025.3531702.

\bibitem{ref64} A. Gonzalez Bermudez, M. Farreras, M. Groshev, J. A. Trujillo, I. de la Bandera, and R. Barco, ``Graph neural networks for O-RAN mobility management: a link prediction approach,'' arXiv:2502.02170, 2025.

\bibitem{ref65} K. Hasan, T. Trappenberg, and I. Haque, ``A generalized transformer-based radio link failure prediction framework in 5G RANs,'' arXiv:2407.05197, 2024.

\bibitem{ref66} M. Santana and K. L. Dias, ``Open RAN-enabled deep learning-assisted mobility management for connected vehicles,'' arXiv:2412.21161, 2024.

\bibitem{ref67} M. V. Narayana, R. Singh, M. R. Deepa, and R. K. Ganti, ``Enabling AI-native mobility in 6G: a real-world dataset for handover, beam management, and timing advance,'' arXiv:2605.12453, 2026.

\bibitem{ref68} S. Wu, M. Alrabeiah, A. Hredzak, C. Chakrabarti, and A. Alkhateeb, ``Deep learning for moving blockage prediction using real millimeter wave measurements,'' in Proc. IEEE International Conference on Communications (ICC), 2022, arXiv:2101.06886.

\bibitem{ref69} C. Lee, H. Cho, S. Song, and J.-M. Chung, ``Prediction-based conditional handover for 5G mm-wave networks: a deep-learning approach,'' \textit{IEEE Vehicular Technology Magazine}, vol. 15, no. 1, pp. 54--62, 2020, doi:10.1109/MVT.2019.2959065.

\bibitem{ref70} I. A. Safiullin, I. P. Ashaev, A. A. Korobkov, A. K. Gaysin, and A. F. Nadeev, ``Reinforcement learning-based handover algorithm for 5G/6G AI-RAN,'' \textit{Inventions}, vol. 11, no. 1, art. 8, 2026, doi:10.3390/inventions11010008.

\bibitem{ref35} D. R. Roberts et al., ``Cross-validation strategies for data with temporal, spatial, hierarchical, or phylogenetic structure,'' Ecography, vol. 40, no. 8, pp. 913--929, 2017, doi:10.1111/ecog.02881.

\bibitem{ref42} D. Hong and S. S. Rappaport, ``Traffic model and performance analysis for cellular mobile radio telephone systems with prioritized and nonprioritized handoff procedures,'' IEEE Transactions on Vehicular Technology, vol. 35, no. 3, pp. 77--92, 1986.

\bibitem{ref43} X. Lin, R. K. Ganti, P. J. Fleming, and J. G. Andrews, ``Towards understanding the fundamentals of mobility in cellular networks,'' IEEE Transactions on Wireless Communications, vol. 12, no. 4, pp. 1686--1698, Apr. 2013, doi:10.1109/TWC.2013.022113.120512.

\bibitem{ref44} S. Sadr and R. S. Adve, ``Handoff rate and coverage analysis in multi-tier heterogeneous networks,'' IEEE Transactions on Wireless Communications, vol. 14, no. 5, pp. 2626--2638, May 2015, doi:10.1109/TWC.2015.2390231.

\bibitem{ref48} R. L. Prentice and L. A. Gloeckler, ``Regression analysis of grouped survival data with application to breast cancer data,'' Biometrics, vol. 34, no. 1, pp. 57--67, 1978, doi:10.2307/2529588.

\bibitem{ref49} P. D. Allison, ``Discrete-time methods for the analysis of event histories,'' Sociological Methodology, vol. 13, pp. 61--98, 1982, doi:10.2307/270718.

\bibitem{ref50} G. Tutz and M. Schmid, \textit{Modeling Discrete Time-to-Event Data}, Springer Series in Statistics, Springer, Cham, 2016, doi:10.1007/978-3-319-31552-2.

\bibitem{ref16} S. Wiegrebe, P. Kopper, R. Sonabend, B. Bischl, and A. Bender, ``Deep learning for survival analysis: a review,'' Artificial Intelligence Review, vol. 57, no. 3, art. 65, 2024, doi:10.1007/s10462-023-10681-3.

\bibitem{ref17} H.-C. Thorsen-Meyer, D. Placido, B. S. Kaas-Hansen, et al., ``Discrete-time survival analysis in the critically ill: a deep learning approach using heterogeneous data,'' npj Digital Medicine, vol. 5, art. 142, 2022, doi:10.1038/s41746-022-00679-6.

\bibitem{ref18} E. J. Cand\`es, L. Lei, and Z. Ren, ``Conformalized survival analysis,'' Journal of the Royal Statistical Society Series B, vol. 85, no. 1, pp. 24--45, 2023, doi:10.1093/jrsssb/qkac004.

\bibitem{ref21} A. N. Angelopoulos, S. Bates, A. Fisch, L. Lei, and T. Schuster, ``Conformal risk control,'' in Proc. International Conference on Learning Representations (ICLR), 2024.

\bibitem{ref22} S. Bates, A. Angelopoulos, L. Lei, J. Malik, and M. I. Jordan, ``Distribution-free, risk-controlling prediction sets,'' Journal of the ACM, vol. 68, no. 6, 2021, doi:10.1145/3478535.

\bibitem{ref19} K. M. Cohen, S. Park, O. Simeone, and S. Shamai (Shitz), ``Calibrating AI models for wireless communications via conformal prediction,'' arXiv:2212.07775, 2022.

\bibitem{ref20} O. Simeone, S. Park, and M. Zecchin, ``Conformal calibration: ensuring the reliability of black-box AI in wireless systems,'' arXiv:2504.09310, 2025.

\bibitem{ref55} K. Rezaei, M. Sattari, J. Aliakbari, T. Svensson, P. Monti, and C. Natalino, ``TRACE-CRC: trajectory-adaptive conformal risk control for multi-step channel state information prediction,'' in \textit{Proc. Symposium on Conformal and Probabilistic Prediction with Applications}, PMLR vol. 329, pp. 984--1009, 2026.

\bibitem{ref51} R. J. Tibshirani, R. Foygel Barber, E. Cand\`es, and A. Ramdas, ``Conformal prediction under covariate shift,'' in \textit{Advances in Neural Information Processing Systems 32 (NeurIPS)}, 2019, pp. 2530--2540.

\bibitem{ref52} R. Foygel Barber, E. J. Cand\`es, A. Ramdas, and R. J. Tibshirani, ``Conformal prediction beyond exchangeability,'' \textit{The Annals of Statistics}, vol. 51, no. 2, pp. 816--845, 2023, doi:10.1214/23-AOS2276.

\bibitem{ref53} I. Gibbs and E. Cand\`es, ``Adaptive conformal inference under distribution shift,'' in \textit{Advances in Neural Information Processing Systems 34 (NeurIPS)}, 2021, pp. 1660--1672.

\bibitem{ref38} M. M. S. Shafi, K. M. Istiaque, S. S. Sowad, and M. T. Kawser, ``Drive-test-based LTE handover dataset for cellular mobility studies in urban Bangladesh,'' Mendeley Data, version 2, 2025, doi:10.17632/n2pvmtyn2j.2.

\bibitem{ref40} F. Chen, M. Ghoshal, E. Nan, P. Dinh, I. Khan, Z. J. Kong, Y. C. Hu, and D. Koutsonikolas, ``A large-scale study of the potential of multi-carrier access in the 5G era,'' in Passive and Active Measurement (PAM), LNCS vol. 15567, 2025, pp. 469--484, doi:10.1007/978-3-031-85960-1\_19. Dataset: github.com/NUWiNS/pam2025-multi-carrier-dataset.

\bibitem{ref37} D. Raca, D. Leahy, C. J. Sreenan, and J. J. Quinlan, ``Beyond throughput, the next generation: a 5G dataset with channel and context metrics,'' in Proc. ACM Multimedia Systems Conference (MMSys), 2020, pp. 303--308, doi:10.1145/3339825.3394938. Dataset: github.com/uccmisl/5Gdataset (GPL-3.0).

\bibitem{ref34} B. Zadrozny and C. Elkan, ``Transforming classifier scores into accurate multiclass probability estimates,'' in Proc. ACM SIGKDD, 2002, pp. 694--699.

\bibitem{ref33} C. Guo, G. Pleiss, Y. Sun, and K. Q. Weinberger, ``On calibration of modern neural networks,'' in Proc. ICML, 2017, pp. 1321--1330.

\bibitem{ref32} G. Ke, Q. Meng, T. Finley, T. Wang, W. Chen, W. Ma, Q. Ye, and T.-Y. Liu, ``LightGBM: a highly efficient gradient boosting decision tree,'' in Advances in Neural Information Processing Systems 30, 2017, pp. 3146--3154.

\end{thebibliography}
